\begin{ledgroupsized}[r]{120mm}
\footnotesize
\pstart
\noindent\textbf{Überlieferung:}
\pend
\end{ledgroupsized}
\begin{ledgroupsized}[r]{114mm}
\footnotesize
\pstart
\parindent -6mm
\makebox[6mm][l]{\textit{L}}%
Aufzeichnungen:
LH~IV~3, 9 Bl.~10.
Ein Blatt~2\textsuperscript{o};
{Wasserzeichen} in der Blattmitte (französisches Papier);
Papiererhaltungsmaßnahmen;
alle Ränder bis auf den rechten beschnitten;
geringfügiger Textverlust durch Papierschaden am rechten Rand.
Die in N.~\ref{Cc2_1381} edierten Aufzeichnungen umfassen Rechnungen und Zeichnungen auf Bl.~10~v\textsuperscript{o};
Bl.~10~r\textsuperscript{o} und der Großteil von Bl.~10~v\textsuperscript{o} überliefern das Konzept \cite{02168}"De veritatibus, de mente, de Deo, de universo" (\textit{LSB} VI,~3 N.~71, S.~507--513) ohne Zusammenhang mit N.~\ref{Cc2_1381}; bei dessen Abfassung hat Leibniz die bestehenden Rechnungen und Zeichnungen teils über- und teils umschrieben.%
\pend
\pstart\noindent
Cc 2, Nr. 1381
\pend
\end{ledgroupsized}
\vspace{5mm}
\begin{ledgroup}
\footnotesize
\pstart
\noindent
\textbf{Datierungsgründe:}
Zur Datierung der Aufzeichnungen in N.~\ref{Cc2_1381} ist zunächst der Umstand maßgebend, dass Leibniz auf demselben Träger auch das \cite{02168}Konzept "De veritatibus" verfasst hat, wobei seine Platznutzung, wie aus der Textverteilung dieser Schicht ersichtlich ist, durch die vorgefundenen Rechnungen und Zeichnungen eingeschränkt war.
Das Konzept hat Leibniz am oberen Rand von 
\href{\Scan{LH_04_03_09b_010r}}%
{Bl.\;10~r\textsuperscript{o}} 
eigenhändig auf den 25. April 1676 datiert (die Hrsg. von \textit{LSB} VI,~3 lesen "15. April").
Dieses Datum bildet daher einen Terminus ante quem für die Entstehung von N.~\ref{Cc2_1381}.
\pend
\pstart
Den Terminus post quem, somit die vorgeschlagene Datierungsspanne, bietet das für Februar bis Juli 1676 häufig belegte Wasserzeichen.
Dieses ist u.\,a. in den
\cite{02169}Auszügen aus 
\indexN{\textls{Mariotte}, Edme, Seigneur de Chazeuil ca. 1620\textendash1684}Mariottes  \textit{Logique}-Handschrift (\textit{LSB} VI, 3 N.\;22, Mai bis Juli 1676)
und den Exzerpten \cite{02170}\textit{Anatomica quaedam ex manuscripto \indexN{\textls{Descartes} (Cartesius, des Cartes), René 1596\textendash1650}Cartesii} mit der editorischen Datierung Februar bis September 1676 (\textit{LSB} VII, 2 N.\;58) belegt.
Dass auch ein späteres Stück aus der Hannoveraner Zeit (N.~\ref{Cc2-1421} von 1679 -- Januar 1687) dasselbe Wasserzeichen im Träger aufweist, erklärt sich in jenem Fall dadurch, dass Leibniz dafür die freigebliebenen Teile eines alten Blattes aus der Pariser Zeit nachträglich verwendete; siehe dazu die Datierungsbegründung von N. \ref{Cc2-1421}.
\pend 
\end{ledgroup}
\clearpage
\selectlanguage{latin}		%
\vspace{8mm}
\pstart
\noindent
\href{\Scan{LH_04_03_09b_010v}}%
{\lbrack10~v\textsuperscript{o}\rbrack} 
\pend
\nobreak
\pstart
\noindent
\bracks{\textit{Im oberen Viertel:}}
\pend
\nobreak
\begin{edfigure}[0.5em][1em]
\hfill
\hspace*{10mm}
\includegraphics[scale=1]{%
gesamttex/edit_VIII-4/images/LH_04_03_09_010_d1_010v.pdf}
\hfill%
\hspace*{9mm}
\includegraphics[scale=1]{%
gesamttex/edit_VIII-4/images/LH_04_03_09_010_d2_010v.pdf}
\hspace*{10mm}
\hfill
\edcaption{\lbrack\textit{Fig.~1}\rbrack 
\hspace*{40mm}%
\lbrack\textit{Fig.~2}\rbrack
\hspace*{20mm}}
\end{edfigure}
\pstart
\textit{AB} \textit{AC} \textit{AD}\\
\pend
\vspace{.5em}
\pstart
\noindent
\bracks{\textit{Rechnungen in der rechten Hälfte:}}
\pend
\vspace{.5em}
\nobreak
\pstart
\noindent
$\begin{array}{r}
33\\33\\\overline{99}\\99\phantom{0}\\\overline{1089}
\end{array}$
\hspace{5mm}
1.\quad $\dfrac{33}{100}$\quad $\dfrac{1089}{100}$
\hspace{5mm}
$\begin{array}{r}
\cancel{1}\phantom{\cancel{0}}\\\cancel{1}\cancel{0}\cancel{0}\\\cancel{3}\cancel{3}
\end{array}f\, 33$
\hspace{5mm}
$\begin{array}{l}
20+\frac{100}{3}\\33\,+\\\overline{53+\frac{1}{3}}
\end{array}$
\pend
\vspace{.3em}
\pstart
\noindent
$\begin{array}{r}
65\phantom{000}\\15000\\\overline{325000}\\65\phantom{5000}\\\overline{975000}
\end{array}$
\hspace{5mm}
$\begin{array}{r}
20\\25\phantom{0}\\\overline{500}
\end{array}$
\hspace{5mm}
$\begin{array}{r}
800\\15\phantom{00}\\\overline{1300}\\\overline{650}
\end{array}$
\hspace{5mm}
300
\pend
\vspace{.5em}
\pstart
\noindent
\bracks{\textit{Reihen, über die ganze Seite verteilt:}}
\pend
\nobreak
\vspace{.5em}
\pstart
\noindent
2 4
\hspace{5mm}
2 4 8 16 \hspace{2mm}
$\begin{array}[t]{r}36\\2\end{array}$
\hspace{5mm}
2 6 18
\pend
\pstart
\noindent
1.2.3.4.5
\hspace{5mm}
$\dfrac{1}{1}$ $\dfrac{1}{2}$ $\dfrac{1}{3}$ $\dfrac{1}{4}$ $\dfrac{1}{5}$
\hspace{5mm}
$\doublefrac{\dfrac{1}{1}}{\dfrac{1}{3}}$ $\doublefrac{\dfrac{1}{2}}{\dfrac{1}{6}}$ $\doublefrac{\dfrac{1}{2}}{\dfrac{1}{4}}$ $\doublefrac{\dfrac{1}{6}}{\dfrac{1}{12}}$
\hspace{5mm}
$\dfrac{1}{1}$ $\dfrac{1}{2}$ $\dfrac{1}{3}$ $\dfrac{1}{4}$
\pend
\pstart
\noindent
$1\cdot 1\cdot 1\cdot $\rule[-5mm]{0pt}{10mm}
\hspace{5mm}
$1+1+1$
\pend
\pstart
\noindent
$\dfrac{1}{2} \bigtimes \dfrac{1}{3}$\rule[-6mm]{0pt}{12mm}\quad $\dfrac{3}{6}$
\hspace{5mm}
$\dfrac{1, \cancel{2}}{\cancel{2}}$ 
\hspace{5mm}
$\dfrac{1}{3}$
\hspace{5mm}
$\dfrac{1}{2}$ $\dfrac{4}{12}$
\pend
\vspace{.3em}
\pstart\noindent
\begin{edtabularc}%
$\displaystyle\efrac{1}{2}$&&&&
$\dfrac{1}{1}$ & $\dfrac{1}{2}$ & $\dfrac{1}{3}$ && $\dfrac{1}{4}$ & \hspace{5.5cm}
\\
$\displaystyle\efrac{\rule{0pt}{1.33em}\overline{\rule{0pt}{.85em}\,3\,}}{6}$&&&&
\;\;24 & \;\;12 & \;\;8 && \;6 &
\\%
$\dfrac{\rule{0pt}{1.2em}4}{24}$&&&&%
$\dfrac{24}{8}$
&
\!$\dfrac{12}{4}$
&
$\dfrac{12}{6}$ & \hspace{-2em}$\sqcap$\hspace{-2em} & $\dfrac{4}{2}$ & \\
\end{edtabularc}
\pend
\vspace{-1em}
\pstart
\noindent
\bracks{\textit{Rechnungen am linken Seitenrand:}}
\pend
\vspace{.5em}
\nobreak
\pstart
$\begin{array}[t]{l}
\phantom{-}3b-1\\\phantom{-}3c-1\\\overline{-3b+1}\\-3c\phantom{+1}\\\underline{+3bc\phantom{+1}}\\3bc-3b=3c+1
\end{array}$
\hspace{5mm}
$3b-1,\fbox{2}$. $\ovalbox{9b^2}-6b+1\sqcap\;\ovalbox{b^2}-2cb+c^2$. $b\sqcap 1\dfrac{b^2+c^2}{}$
\pend
\vspace{.5em}
\pstart
\noindent
\bracks{\textit{Rechnungen in der oberen Hälfte:}}
\pend
\vspace{.5em}
\nobreak
\pstart
$\dfrac{b^2-e^2}{b-e}\sqcap\;b+e$
\hspace{5mm}
$\begin{array}{r}
b+e\\+b-e\\\overline{+b^2-eb-e^2}\\+eb\phantom{-be^2}\\\overline{+b^2\phantom{-eb}-e^2}
\end{array}$
\pend
\pstart
$a\ b\ y$
\hspace{5mm}
$\dfrac{a.}{b.}$ $\dfrac{b}{y}$ $\dfrac{y}{\cancel{b}}\sqcap\dfrac{bb}{a}$
$y\sqcap\dfrac{b^2}{a}$
\hspace{5mm}
$\dfrac{\cancel{yb}}{\cancel{b^2}}$
\pend
\clearpage
\pstart
\noindent
\bracks{\textit{Zeichnungen, Text und Rechnungen im unteren Viertel:}}
\pend
\nobreak
\begin{edfigure}[1em][1em]
\hfill
\includegraphics[scale=1]{%
gesamttex/edit_VIII-4/images/LH_04_03_09_010_d3_010v.pdf}
\hfill%
\includegraphics[scale=1]{%
gesamttex/edit_VIII-4/images/LH_04_03_09_010_d4_010v.pdf}
\hfill
\edcaption{%
\hspace*{5mm}%
\lbrack\textit{Fig.~3}\rbrack 
\hspace*{42mm}%
\lbrack\textit{Fig.~4}\rbrack}
\end{edfigure}
\pstart
\selectlanguage{french}
{
\begin{edtabularl}%
&lieue&\hspace{5.5cm}\\
1500&longueur d'un canal&\\
$b\sqcap AB\sqcap 40$&pieds largeur superieure&\\
$g\sqcap DE,\sqcap FG\sqcap 20$&profondeur&\\
$e\sqcap EH\sqcap 25$&pieds largeur inferieure&\\
\end{edtabularl}%
}%
\pend
\vspace{.5em}
\pstart
\noindent
$AB\sqcap 40$\hspace{5mm} $\dfrac{AB\sqcap b}{EH\sqcap e}\sqcap\dfrac{FC\sqcap g+y}{GC\sqcap y}$
\pend
\pstart
\noindent
$\dfrac{b}{e}\sqcap\dfrac{g+y}{y}$.
\hspace{5mm}
$yb\sqcap eg+ey$.
\hspace{5mm}
$yb-ey\sqcap eg$.
\hspace{5mm}
$y\sqcap \dfrac{eg}{b-e}\sqcap\dfrac{25, 20}{15}\sqcap\dfrac{30,5}{3}$
\pend
\vspace{.25em}
\pstart
\noindent
$GC\sqcap y\sqcap\dfrac{100}{3}$ 
\hspace{5mm}
$FC\sqcap g+\dfrac{eg}{b-e}\sqcap\dfrac{bg}{b-e}$
\hspace{5mm}
$ABC\sqcap\dfrac{gb}{2}+\dfrac{egb}{2,b-e}\sqcap\dfrac{gb^2}{2,b-e}$
\pend
\vspace{.35em}
\pstart
\noindent
$\dfrac{e^2g}{2,b-e}$
\hspace{5mm}
$\dfrac{gb}{\cancel{2}}$
\hspace{5mm}
$\dfrac{gb^2\ovalbox{-gbe+gbe}-e^2g}{b-e}$
\hspace{5mm}
$\dfrac{gb^2-ge^2}{b-e}\sqcap\ovalbox{\dfrac{gb+ge}{2}}$
\pend
